\section{Introduction}
\label{sec:intro}

Renormalization of lattice operators
is a necessary ingredient to obtain many physical
results from numerical simulations, such as meson decay constants,
form factors, structure functions, mixing amplitudes, etc.
In this paper we study the
renormalization of composite operators in lattice $QCD$.
We propose a method which avoids completely
 the use of lattice perturbation theory and allows a non-perturbative
determination
of the renormalization constants of any composite operator.
For illustrative purposes, we consider some specific applications
to matrix elements of two-fermion operators.
The approach is particularly useful in those cases,
such as the scalar or pseudoscalar densities,
where it is not possible to use the chiral Ward identities to determine
the renormalization constants non-perturbatively\footnote{In the following
we will denote the method to determine the renormalization constants, using
the Ward identities, as the Ward identity method.}  \cite{bo}--\cite{wi}.
Moreover,
our proposal can be applied to many other cases,
such as   the renormalization of
the four-fermion operators of the effective weak
Hamiltonian, and to heavy-light currents in the Heavy Quark Effective
Theory. Preliminary results have been presented in ref.\ \cite{dallas}.
A similar attempt, limited to the divergent part of the $\Delta I=1/2$
penguin-operator, can be found in ref.\ \cite{ber1}.

Renormalized lattice operators must
correspond, in the limit of
infinite lattice cut-off, to finite chirally covariant operators, which
obey the same renormalization conditions as those in the continuum.
 With just a few exceptions, lattice perturbation theory is
 used to evaluate
the  renormalization constants of lattice operators.
The problem of mixing with lower dimensional
operators requires however, a non-perturbative
subtraction of all power divergences \cite{bo},\cite{te}--\cite{martsukuba}.
Apart from this special case,
the use of perturbation theory, in the computation of
multiplicative  constants and dimensionless mixing coefficients,
is well justified,
provided that the lattice spacing $a$ is sufficiently small,
i.e.\ $ a^{-1} \gg \Lambda_{QCD}$.
However, in those cases where the results from perturbation theory,
obtained using the bare lattice
coupling constant $\asl=g_{0}^{2}/4 \pi$ as an
expansion parameter, can be checked
using non-perturbative methods, there is generally a significant
discrepancy.
Several solutions to this problem have been proposed so far:
\begin{itemize}
\item
In order to improve the convergence
of the perturbative series,
the authors of ref.\ \cite{lepage} have  proposed a
``tadpole improved'' perturbation theory.
Nevertheless, ignorance of higher order contributions still  represents an
important source
of uncertainty in the extraction of physical results.
\item
In some limited cases, corresponding to finite
operators, namely the vector and axial vector currents,
and  the ratio of pseudoscalar and scalar densities, a fully
non-perturbative determination of the renormalization constants can be
obtained with the use of chiral Ward identities \cite{bo}--\cite{wi}.
\item
One may fix  non-perturbative renormalization conditions directly on
hadronic matrix elements. This procedure was used in ref.\ \cite{gav1},
to subtract the divergent part of the $\Delta I=1/2$ operator.
The price is that one has to sacrifice
 a physical prediction, for any subtraction imposed
on hadronic states. Thus, when several
renormalization
conditions are necessary,
 one  looses much predictive power \cite{bo,te}.
\end{itemize}

{\itshape
Our proposal is to impose
re\-nor\-ma\-li\-za\-tion con\-di\-tions
non-per\-tur\-ba\-ti\-vely,  di\-rectly on quark and
gluon Green functions, in a fixed gauge, with given off-shell
ex\-ter\-nal sta\-tes, with large vir\-tua\-li\-ties.}

The method consists in  mimicking  what is usually done in perturbation
theory.
One fixes the renormalization conditions of a
certain operator by imposing that suitable Green functions, computed
between external off-shell quark and gluon states, in a fixed gauge,
coincide with their tree level value. For example, if we consider
the generic two-quark operator $O_\Gamma= \bar \psi \Gamma \psi$,
we may impose the condition
\beq Z_\Gamma \langle p \vert O_\Gamma \vert p \rangle\vert_{p^2=-\mu^2}=
\langle p \vert O_\Gamma \vert p \rangle_0, \eeq
where $\Gamma$ is one of the Dirac matrices and $
\langle p \vert O_\Gamma \vert p \rangle_0$ is the tree level matrix element.
The extension to more complicated cases, including four-fermion operators
and operator mixing, is straightforward.
This procedure defines the same renormalized operators, i.e.\ the same
Wilson coefficient functions, in all  regularization schemes,
 provided they are expressed in terms of the same renormalized
coupling constant. However, the coefficient functions now depend on
the external states and on the gauge, which must be specified.

In principle this scheme completely solves
the problem of  large corrections in lattice perturbation theory,
 which are automatically
included in the renormalization constants.
Matching of  lattice operators to
the corresponding ones, renormalized in a
continuum renormalization scheme (for definiteness, the $\overline{MS}$
scheme),
then requires continuum perturbation theory only.
When  next-to-leading corrections are known, the error
on the physical matrix elements is of  $O((\as)^2)$, in the continuum
perturbative expansion, which is expected to have a better convergence.
A continuum perturbative calculation of the matching conditions
 is a common step in  all approaches, standard or
 ``tadpole improved" perturbation theory and non-perturbative
renormalization.
Our proposal  can be applied to any composite operator,
unlike the  Ward identity method, which
 is limited to a few, albeit very important, cases\footnote
{It is still true that the statistical
accuracy is in general slightly better with the Ward identities
method, in those cases
where it can be applied, see sec.\ \ref{sec:nonper}.}.

The feasibility of using Green functions computed between
 quark states,
in a fixed gauge, was already discussed in ref.\ \cite{wi}, in the context
of non-perturbative renormalization, using Ward identities. Indeed, in the case
of the axial current, it was shown
that the signal was rather good and the result was in agreement with the
corresponding hadronic gauge invariant determination.

As far as the renormalization conditions are concerned,
our proposal is expected to work, whenever it is  possible to
fix the virtuality of the external states $\mu$ and to satisfy
the condition $ \Lambda_{QCD} \ll \mu \ll 1/a $, in order
to keep under control both non-perturbative and discretization effects.
We stress that this  requirement is common to all methods.
The results of our numerical investigation
are encouraging: they suggest that there does exist a ``window" in $\mu$,
where the method can be applied,
at values of the lattice coupling constant used in
current lattice simulations.

The plan for the remainder of the paper is as follows. In the next
section the basic formulae necessary to define our non-perturbative
method are presented, using as an example logarithmically divergent
two-quark operators, such as the scalar and pseudoscalar densities. The
discussion is generalised to all composite operators in section 3, and
in the following section the implementation of the method in numerical
simulations for two-quark operators is discussed. In section 5 we give
some details about the perturbative evaluation of the renormalisation
constants on lattices of finite volume. The values of the
renormalisation constants obtained in a numerical simulation are
presented in section 6, and are compared to the results from
perturbation theory. The final section contains our conclusions, a
summary of the numerical results and a discussion of possible further
applications of the ideas presented in this paper.
